\begin{afterword}
\summaryhead

The post argues that a belief should move by small steps as evidence comes in, down as well as up. A theory that predicts probabilistically will meet some contrary evidence even when it is true: a coin that lands heads 95\% of the time still lands tails one time in twenty. So one contrary observation calls for a small downward shift, not a defence that argues it away and not a rejection.

The author claims that people reason in yes-or-no terms, as if a true theory could have no failures and a false one no successes, and that they want debates in which their side concedes nothing. To the objection that concessions lose debates, the post answers that rationality is for deciding which side to join. By conservation of expected evidence, you must expect as much downward revision as upward. When contrary evidence keeps coming, you have already given the belief up, and should celebrate.

\responsehead

The post's lesson is correct and useful. A theory that assigns probabilities to outcomes should expect some outcomes against it, and one of them is not a refutation. The coin case shows this well: nineteen heads for each tail is exactly what the 95\% hypothesis predicts. The post also states the conservation rule more exactly than the post it cites. The average of your possible later beliefs equals your present one, which is true, and it avoids that post's slip about expecting to stay ``exactly as confident.''

The post does not say how large a step should be. The post tells the reader to ``think quantitatively'' and to ``shift downward a little,'' and it gives no way to find the quantity. How far one observation should move you depends on what the rival hypothesis predicts, and the post never names a rival. With a fair coin as the rival, a single tail divides the odds for the 95\% hypothesis by ten, which is hardly ``a little.'' The post's own restatement of the conservation rule, two paragraphs later, holds that rare contrary evidence must move belief a lot. So the comforting advice holds for weak evidence, which the post calls an ``iota,'' and a reader cannot tell from the post which evidence is weak.

The post describes the opposing view without evidence. The post claims that it is ``widely believed'' that a true theory can have no failures, and ``this is how humans have argued for ages and ages.'' No one who argues this way is named or quoted. The post may well be right about how people argue, but it offers only the description.

The objection from debating is answered by changing the goal. The objector says that a debater who concedes a point will lose. The post does not dispute this; it answers that rationality is for deciding which side to join. That is a fair reply for someone deciding what to believe, and it leaves the debater's question unanswered.

In short: a correct rule, that contrary evidence should lower a belief without overturning it, given with no way to tell how far any one observation should move it.
\end{afterword}
